\section{Sucesión de Mayer-Vietoris}

Una familia de grupos abelianos $\{A_n:n\in \N\}$ y morfismos $\alpha_n: A_n\to A_{n-1}$ constituye una $\textbf{sucesión exacta}$ si para todo $n\in\N$ se cumple $\Ker(\alpha_{n})=\im(\alpha_{n+1})$.
\begin{obs}

La inclusión $\im(\alpha_{n+1})\subset\Ker(\alpha_{n})$ nos dice que $\alpha_{n+1}\circ\alpha_{n}\equiv 0$, es decir, los $\alpha_n$ forman un complejo de cadenas. 

La otra inclusión, $\im(\alpha_{n+1})\supset\Ker(\alpha_{n})$, nos dice que los grupos de homología asociados a este complejo de cadenas son triviales.

En general, podemos interpretar la sucesión de grupos de homología asociados a un complejo de cadenas como un indicador de qué tan cerca está el complejo de ser una sucesión exacta.
\end{obs}
Llamaremos \textbf{sucesión exacta corta} a una sucesión exacta de la forma $$0\rightarrow A\xrightarrow{f} B\xrightarrow{g} C\rightarrow 0$$
\begin{obs}
	Una sucesión de esta forma es exacta si y solo si $f$ es inyectiva, $g$ es sobreyectiva y $\Ker(g)=\im(f)$. Por lo tanto, $g$ induce un isomorfismo entre $C$ y $B/\im(f)$.
\end{obs}
En general, una \textbf{sucesión exacta larga} es cualquier sucesión exacta con más términos que una sucesión exacta corta. En este capítulo, usaremos el término para referirnos específicamente a sucesiones exactas con una cantidad infinita numerable de elementos.

El foco de este capítulo es describir la existencia de una sucesión exacta larga particular, conocida como la \textbf{sucesión de Mayer-Vietoris}.

La existencia de la sucesión de Mayer-Vietoris para espacios topológicos va a surgir del siguiente hecho algebraico:

\begin{lema}{Lema de la serpiente}

	Sean $A$, $B$ y $C$ espacios tal que para cada $n\in\N$ existen morfismos $f_n\colon C_n(A)\to C_n(B)$ y $g_n\colon C_n(B)\to C_n(C)$ para los que la sucesión $$0\rightarrow C_n(A)\xrightarrow{f_n} C_n(B)\xrightarrow{g_n} C_n(C)\rightarrow 0$$ es una sucesión exacta corta entre complejos de cadena. 
	
	Entonces, para cada $n\in\Z$ existe un morfismo $$\delta_* :H_{n+1}(C)\rightarrow H_n(A)$$ tal que la siguiente sucesión es exacta:

$$\begin{tikzcd}[arrow style=math font,cells={nodes={text height=2ex,text depth=0.75ex}}]
       \cdots \arrow[overlay, r, "f_{n+1}^*"] & H_{n+1}(B) \arrow[overlay, r, "g_{n+1}^*"] \arrow[overlay, draw=none]{d}[name=Y, shape=coordinate]{} & H_{n+1}(C) \arrow[overlay, curarrow=Y]{dll}{} \\
       H_{n}(A) \arrow[overlay, r, "f_n^*"] & H_{n}(B) \arrow[overlay, r, "g_n^*"] \arrow[overlay, draw=none]{d}[name=Z,shape=coordinate]{} & H_{n}(C) \arrow[overlay, curarrow=Z]{dll}{} \\
       H_{n-1}(A) \arrow[overlay, r, "f_{n-1}^*"] & H_{n-1}(B) \arrow[overlay, r, "g_{n+1}^*"] & \cdots
\end{tikzcd}$$
\end{lema}

Dentro de un espacio topológico consideramos dos subespacios $A$ y $B$, y sea $X=\mathring{A}\cup\mathring{B}$. Definimos $C_n(A+B)$ el complejo de cadenas dado por la suma formal de elementos de $C_n(A)$ y $C_n(B)$. El mapa borde $\partial:C_{n}(X)\to C_{n-1}(X)$ manda elementos de $C_n(A+B)$ en elementos de $C_{n-1}(A+B)$, entonces los $C_n(A+B)$ forman un complejo de cadenas.

\[\begin{tikzcd}
	\dots & {C_n(A+B)} & {C_{n-1}(A+B)} & \dots \\
	\dots & {C_n(X)} & {C_{n-1}(X)} & \dots
	\arrow[from=1-1, to=1-2]
	\arrow[from=1-2, to=1-3]
	\arrow["i", hook, from=1-2, to=2-2]
	\arrow[from=1-3, to=1-4]
	\arrow["i", hook, from=1-3, to=2-3]
	\arrow[from=2-1, to=2-2]
	\arrow[from=2-2, to=2-3]
	\arrow[from=2-3, to=2-4]
\end{tikzcd}\]

$\textbf{Afirmación:}$ Las inclusiones $C_n(A+B)\hookrightarrow C_n(X)$ inducen isomorfismos en los grupos de homología. 

$$H_n(X)\simeq H_n(A+B)$$

\comm{me falta agregar esta prueba}

Definiendo $\psi(x)=(x,-x)$, $\phi(x,y)=x+y$ consideramos el siguiente diagrama:

$$0\rightarrow C_n(A\cap B)\xrightarrow{\psi}C_n(A)\oplus C_n(B)\xrightarrow{\phi}C_n(A+B)\rightarrow 0$$

Esta va a resultar una sucesión exacta corta, vamos a verificarlo.

Veamos primero que $\Ker(\phi)=\im(\psi)$: 

Si se cumple $(x,y)\in\Ker(\phi)$ para dos elementos $x\in C_n(A)$, $y\in C_n(B)$, entonces se cumple $x+y=0$. Esto implica que $x=-y$ como cadenas en $X$, y por lo tanto $x, y\in C_n(A\cap B)$. Tenemos entonces $\psi(x)=(x,-x)=(x,y)\in C_n(A)\oplus C_n(B)$, es decir, $(x,y)\in\im(\psi)$. Con esto probamos $\Ker(\phi)\subset \im(\psi)$. La otra inclusión es más directa, dado $(x,-x)\in\im(\psi)$ se tiene $x\in C_n(A\cap B)$, entonces $\phi(x,-x)=x-x=0$

% \duda{me parece medio chafi como esta definido el complejo de A+B}

Además de eso, dado $x\in C_n(A\cap B)$ tal que $\psi(x)=0$, se tiene $(x,-x)=(0,0)$ en $C_n(A)\oplus C_n(B)$, por lo que necesariamente $x=0$ y por lo tanto $\Ker(\psi)=0$.

Lo último que basta verificar es $\im(\phi)=C_n(A+B)$. Efectivamente, dado $x\in C_n(A+B)$, $x$ es una suma formal de $n$-cadenas en $A$ y $B$

% \comm{en algún lugar de acá mencionar eso de que la sucesión de grupos de homología es una forma de expresar qué tan cerca de ser exacta está la sucesión de grupos abelianos}

% \begin{lema}
    
%     Sea $X$ un espacio topológico y $\mathcal{U}=\{U_j\}$ una familia de subespacios tales que $\{\mathring{U_j}\}$ es un cubrimiento de $X$. 

%     Definimos como $C_n^{\mathcal{U}}(X)$ al subgrupo de $C_n(X)$ dado por las cadenas $\sum_i n_i\sigma_i$

% \end{lema}


